% !TeX program = XeLaTeX
% !TeX root = ../../shloka.tex
% --meta--
% deity: Kubera
% type: Shatanama Stotram
% chandas: Anushtup
% source: a Tamil-script Kubera pūjā booklet, pp. 16-17 (Vaidika Pandit Sri V Hariharan Collection,
%   digitised by Muthulakshmi Research Academy, CC-0), transliterated 2026-11-09
% --end-meta--
\sect{कुबेराष्टोत्तरशतनामस्तोत्रम्}
\centerline{\small{(मूलम्—कुबेर-पूजा-पुस्तिका (तमिऴ्-लिपिः), पृ. १६–१७)}}

\twolineshloka
{कुबेरो धनदः श्रीमान् यक्षेशो गुह्यकेश्वरः}
{निधीशः शङ्करसखो महालक्ष्मीनिवासभूः} %॥१॥

\twolineshloka
{महापद्मनिधीशश्च पूर्णः पद्मनिधीश्वरः}
{शङ्खाख्यनिधिनाथश्च मकराख्यनिधिप्रियः} %॥२॥

\twolineshloka
{सुकच्छपनिधीशश्च मुकुन्दनिधिनायकः}
{कुन्दाख्यनिधिनाथश्च नीलनिध्यधिपो महान्} %॥३॥

\twolineshloka
{वरनिध्यधिपः पूज्यो लक्ष्मीसाम्राज्यदायकः}
{तथैवेलपिलापत्यः कोशाधीशः कुलोचितः} %॥४॥

\twolineshloka
{अश्वारूढो विश्ववन्द्यो विशेषज्ञो विशारदः}
{नलकूबरतातश्च मणिग्रीवपिता तथा} %॥५॥

\twolineshloka
{गूढमन्त्रो वैश्रवणश्चित्रलेखामनःप्रियः}
{एकपिङ्गोऽलकाधीशः पौलस्त्यो नरवाहनः} %॥६॥

\twolineshloka
{कैलासशैलनिलयो राज्यदो रावणाग्रजः}
{चित्रचैत्ररथोद्यानविहारसुकुतूहलः} %॥७॥

\twolineshloka
{महोत्साहो महाप्राज्ञः सदापुष्पकवाहनः}
{सार्वभौमोऽङ्गनाथश्च सोमः सौम्यदिगीश्वरः} %॥८॥

\twolineshloka
{पुण्यात्मा पुरुहूतश्रीः सर्वपुण्यजनेश्वरः}
{नित्यकीर्तिर्निधिवेत्ता लङ्काप्राक्तननायकः} %॥९॥

\twolineshloka
{यक्षः परमशान्तात्मा यक्षराड् यक्षिणीवृतः}
{किन्नरेशः किम्पुरुषनाथः खड्गायुधो वशी} %॥१०॥

\twolineshloka
{ईशानदक्षपार्श्वस्थो वायुवामसमाश्रयः}
{धर्ममार्गैकनिरतो धर्मसम्मुखसंस्थितः} %॥११॥

\twolineshloka
{नित्येश्वरो धनाध्यक्षो ह्यष्टलक्ष्म्याश्रितालयः}
{मनुष्यधर्मा सद्वृत्तः कोशलक्ष्मीसमाश्रितः} %॥१२॥

\twolineshloka
{धनलक्ष्मीनित्यवासो धान्यलक्ष्मीनिवासभूः}
{अश्वलक्ष्मीसदावासो गजलक्ष्मीस्थिरालयः} %॥१३॥

\twolineshloka
{राज्यलक्ष्मीजन्मगेहो धैर्यलक्ष्मीकृपाश्रयः}
{अखण्डैश्वर्यसंयुक्तो नित्यानन्दः सुखाश्रयः} %॥१४॥

\twolineshloka
{नित्यतृप्तो निधित्राता निराशो निरुपद्रवः}
{नित्यकामो निराकाङ्क्षो निरुपाधिकवासभूः} %॥१५॥

\twolineshloka
{शान्तः सर्वगुणोपेतः सर्वज्ञः सर्वसम्मतः}
{शर्वाणीकरुणापात्रः सदानन्दकृपालयः} %॥१६॥

\twolineshloka
{गन्धर्वकुलसंसेव्यः सौगन्धिकसुमप्रियः}
{सुवर्णनगरीवासो निधिपीठसमाश्रयः} %॥१७॥

\twolineshloka
{महामेरूत्तरस्थायी महर्षिगणसंस्तुतः}
{तुष्टः शूर्पणखाज्येष्ठः शिवपूजारतोऽनघः} %॥१८॥

\threelineshloka
{राजयोगसमायुक्तो राजशेखरपूजकः}
{नाम्नामष्टोत्तरशतं कुबेरस्य महात्मनः}
{नित्यं पठति यो भक्त्या पुत्रपौत्रसुखी भवेत्} %॥१९॥

{॥इति श्री-कुबेराष्टोत्तरशतनामस्तोत्रं सम्पूर्णम्॥}
